\documentclass[../main]{subfiles}
\ifSubfilesClassLoaded{\setcounter{chapter}{4}}{}
\begin{document}

\chapter{GameObject, Component, Prefab, dan Scene}

\begin{subcpmk}
  \item Sub-CPMK 3.2: Membuat Prefab dan mengelola Scene serta Build Settings
\end{subcpmk}

\section{MonoBehaviour dan Lifecycle}

GameObject dan Component diatur lewat script. Bab ini membahas dasar scripting C\# dan lifecycle \class{MonoBehaviour}. Semua script Unity mewarisi \class{MonoBehaviour} \cite{unityScriptingAPI}. Class ini menyediakan method yang dipanggil engine secara otomatis pada waktu tertentu. Urutan pemanggilan penting untuk memahami kapan kode Anda dijalankan.

\begin{table}[htbp]
\centering
\begin{tabular}{|l|p{6.5cm}|}
\hline
\textbf{Method} & \textbf{Fungsi} \\
\hline
Awake() & Dipanggil sekali saat object dibuat (sebelum Start). Gunakan untuk inisialisasi yang tidak bergantung object lain \\
\hline
Start() & Dipanggil sekali sebelum frame pertama Update. Gunakan untuk setup yang butuh object lain sudah ada \\
\hline
Update() & Dipanggil setiap frame (frekuensi bergantung FPS). Untuk input, logika non-fisika \\
\hline
FixedUpdate() & Dipanggil pada interval tetap (default 50 Hz, timestep 0.02 s). Untuk fisika, Rigidbody, AddForce \\
\hline
LateUpdate() & Dipanggil setelah semua Update. Untuk kamera follow (setelah player bergerak) \\
\hline
\end{tabular}
\caption{Lifecycle Method MonoBehaviour}
\end{table}

\begin{catatan}
\textbf{Fixed Timestep}: Physics engine Unity berjalan pada interval tetap (default 0.02 detik = 50 Hz). \method{FixedUpdate} dipanggil bersamaan dengan physics step sehingga fisika konsisten terlepas dari FPS. \textbf{Golden rule}: baca input di \method{Update}, terapkan gaya di \method{FixedUpdate}. Jangan panggil AddForce di Update---FPS tidak stabil akan membuat gerakan aneh.
\end{catatan}


\section{Component, Prefab, dan Scene}

\begin{learningoutcomes}
\item Menambahkan dan mengkonfigurasi komponen Unity
\item Membuat dan menggunakan Prefab
\item Mengelola Scene dan Build Settings
\end{learningoutcomes}

\subsection{Menambahkan Komponen}

Rigidbody 2D untuk fisika (Gravity Scale = 0 untuk space). Box Collider 2D atau Polygon Collider 2D untuk deteksi tabrakan. Add Component di Inspector.

\subsection{Prefab}

\begin{definisi}[Prefab]
Prefab adalah template GameObject yang dapat digunakan kembali dan diinstansiasi berkali-kali \cite{UnityManual2024}.
\end{definisi}

Membuat Prefab: Buat GameObject (Bullet, Asteroid), tambahkan komponen, drag ke folder Prefabs, hapus dari Hierarchy. Gunakan untuk objek yang di-spawn berulang.

\subsection{Scene Management}

Simpan scene (Ctrl+S) di folder Scenes. File $\rightarrow$ New Scene untuk scene baru. File $\rightarrow$ Build Settings: tarik scene ke Scenes In Build, atur urutan (MainMenu index 0).


\begin{aktivitas}
  \item \textbf{Space Survivor}: Buat project 2D, buat Player (Sprite Triangle), tambahkan Rigidbody 2D dan Collider.
  \item \textbf{Prefab}: Buat Prefab Bullet dan Asteroid, simpan di folder Prefabs.
  \item \textbf{Scene}: Buat scene MainMenu dan Gameplay, tambahkan ke Build Settings.
\end{aktivitas}

\begin{latihan}
  \item Jelaskan perbedaan GameObject dan Component!
  \item Apa fungsi Transform? Jelaskan parent-child relationship!
  \item Kapan menggunakan Prefab?
  \item Bagaimana menambahkan scene ke Build Settings?
\end{latihan}

\begin{asesmen}
\textbf{Instrumen}: Membuat project Space Survivor dengan Player, Prefab Bullet/Asteroid, dan 2 scene.
\textbf{Rubrik}: Lihat Lampiran A
\end{asesmen}

\begin{checklist}
  \item Saya dapat membuat GameObject dan mengatur Hierarchy
  \item Saya dapat menggunakan Transform
  \item Saya dapat menambahkan Rigidbody dan Collider
  \item Saya dapat membuat Prefab
  \item Saya dapat mengelola Scene
\end{checklist}

\begin{rangkuman}
GameObject adalah objek dasar Unity. Transform mengatur posisi, rotasi, skala. Prefab adalah template reusable. Scene adalah container untuk level. Komponen Rigidbody dan Collider untuk fisika dan tabrakan.
\end{rangkuman}

\ifSubfilesClassLoaded{
  \renewcommand{\bibname}{Daftar Pustaka}
  \bibliographystyle{plain}
  \bibliography{../references}
}{}
\end{document}
